% Generated draft: requires mathematical and source review.
\documentclass[11pt]{article}
\usepackage[margin=25mm]{geometry}
\usepackage{fontspec}
\setmainfont{Latin Modern Roman}
\usepackage{amsmath,amssymb,mathtools,unicode-math}
\setmathfont{Latin Modern Math}
\usepackage{enumitem,needspace,xurl,hyperref}
\setlength{\parindent}{0pt}
\setlength{\parskip}{5pt}
\setlength{\emergencystretch}{3em}
\newcommand{\N}{\mathbb{N}}
\newcommand{\Z}{\mathbb{Z}}
\newcommand{\Q}{\mathbb{Q}}
\newcommand{\R}{\mathbb{R}}
\newcommand{\C}{\mathbb{C}}
\newcommand{\F}{\mathbb{F}}
\newcommand{\A}{\mathbb{A}}
\newcommand{\PP}{\mathbb{P}}
\newcommand{\E}{\mathbb{E}}
\newcommand{\eps}{\varepsilon}
\newcommand{\dd}{\,\mathrm{d}}
\begin{document}
\section*{UnsolvedMath ID: 30004995}
\textbf{Source problem code:} catalogue entry
\subsection{Definitions and mathematical statement}
The supplied record asks for a Borel analogue of an incompletely reproduced continuous theorem. The following gives a conditional formulation in a customary symbolic-dynamics setting; the parameter ranges and the meaning of hyper-aperiodicity require confirmation from the full theorem.

Suppose that \(\Gamma\) is a countable discrete group and that \(k\) and \(q\) are nonempty finite discrete alphabets. These ambient hypotheses are not explicitly supplied in the record. Give \(q^\Gamma\) its product topology and the left shift action
\[
(g\cdot y)(h)=y(g^{-1}h).
\]
Write \(\operatorname{Stab}_\Gamma(y)=\{g\in\Gamma:g\cdot y=y\}\) and
\[
\operatorname{Free}(q^\Gamma)=\{y\in q^\Gamma:\operatorname{Stab}_\Gamma(y)=\{e\}\},
\]
with analogous notation for \(k^\Gamma\). A subshift of finite type is a set \(S\subseteq q^\Gamma\) for which there are a finite set \(D\subseteq\Gamma\) and a set \(P\subseteq q^D\) such that
\[
S=\{y\in q^\Gamma:(g\cdot y)|_D\in P\text{ for every }g\in\Gamma\}.
\]
For an action \(\Gamma\curvearrowright X\) and a map \(c:X\to q\), define its equivariant coding map by
\[
\pi_c(x)(g)=c(g^{-1}\cdot x).
\]
An \(S\)-coloring is a map \(c\) satisfying \(\pi_c(X)\subseteq S\). For the supplied \(S\)-coloring \(f:\operatorname{Free}(k^\Gamma)\to q\), the stated hypothesis on its image is
\[
\bigcap_{x\in\operatorname{Free}(k^\Gamma)}\operatorname{Stab}_\Gamma(\pi_f(x))=\{e\}.
\]
This is a common-stabilizer condition on the entire image; it does not assert that every image point has trivial stabilizer.

A customary definition calls \(y\in q^\Gamma\) hyper-aperiodic when
\[
\overline{\Gamma\cdot y}\subseteq\operatorname{Free}(q^\Gamma).
\]
Under the pointwise interpretation, a coloring \(c\) is hyper-aperiodic when every \(\pi_c(x)\) has this property. The record does not specify whether it instead requires the stronger uniform condition \(\overline{\pi_c(X)}\subseteq\operatorname{Free}(q^\Gamma)\).

\textbf{Problem, with unresolved qualifications.} Determine suitable additional hypotheses \(H\) on the group, alphabets, subshift, or Borel map \(f\), under which the stated common-stabilizer condition implies that every free Borel action \(\Gamma\curvearrowright X\) on a standard Borel space admits a Borel \(S\)-coloring \(c:X\to q\) satisfying
\[
\forall x\in X,\qquad
\overline{\Gamma\cdot\pi_c(x)}\subseteq\operatorname{Free}(q^\Gamma).
\]
If the source intends uniform hyper-aperiodicity, replace this last requirement by \(\overline{\pi_c(X)}\subseteq\operatorname{Free}(q^\Gamma)\). Here freeness means that \(g\cdot x=x\) implies \(g=e\). The use of standard Borel spaces is part of the conditional framework and is not confirmed by the supplied excerpt.

The source explicitly says that merely replacing continuity of \(f\) by Borel measurability in its Theorem 3 is insufficient, and that even the correct formulation is unclear. Theorem 3 and the concrete test question announced at the end of the excerpt are absent. Consequently this is a qualified formulation of a request for additional hypotheses, not a fully recovered source theorem or an assertion that the unrestricted Borel implication holds.
\subsection{Short English statement}
Find additional assumptions under which a Borel symbolic coloring with trivial common stabilizer of its coded image yields a Borel hyper-aperiodic coloring for every free Borel action. The source says that a direct replacement of continuity by Borel measurability fails. Its full theorem, parameter restrictions, and intended hyper-aperiodicity condition must be checked before this conditional formulation can be finalized.
\subsection{Sources}
\begin{itemize}
\item Local problem registry: permanent UnsolvedMath ID \(30004995\), source problem code \texttt{catalogue entry}, title \emph{Borel Extensions of Hyper-Aperiodic Coloring Theorems}; collection: Oberwolfach Reports Open Problems. The supplied proposed year is 2022; no proposer is identified.
\item Supplied source citation: \emph{Groups and Dynamics: Topology, Measure, and Borel Structure (2023)}, Oberwolfach Reports, pp.\ 107--169, DOI \texttt{10.4171/owr/2022/3}. \url{https://doi.org/10.4171/owr/2022/3}. The literature-source label is \emph{Oberwolfach Report 3/2022}. The differing year labels are retained without resolution.
\item Supplied literature reference: Bernshteyn, \emph{Equivariant maps to subshifts whose points have small stabilizers}. \url{https://arxiv.org/abs/2106.09673}. Its inclusion does not identify Bernshteyn as the proposer of this registry problem.
\item Literature assessment dated exactly \texttt{2026-08-22}: the supplied curation labels the problem open and states that additional hypotheses are necessary for a Borel analogue of the continuous theorem. This is supplied dated curation, not an independently verified assessment of the literature or current status.
\item The registry labels the statement \texttt{corrected\_verified} and reports reconstruction and checking against the official report. These are supplied provenance claims, not verification performed for this draft. The excerpt attributes the proofs to combinatorial methods, particularly a continuous Lov\'asz Local Lemma cited there as reference [1]; the bibliographic identity of [1] is not supplied.
\end{itemize}
\end{document}
